\documentclass[11pt,a4paper]{article}
\usepackage[T1]{fontenc}\usepackage{lmodern,amsmath,amssymb,amsthm,mathtools,booktabs,microtype}
\usepackage[margin=27mm]{geometry}\usepackage[colorlinks=true,linkcolor=blue,urlcolor=blue,citecolor=blue]{hyperref}
\newtheorem{theorem}{Theorem}[section]\newtheorem{lemma}[theorem]{Lemma}
\newcommand{\Q}{\mathbb Q}\newcommand{\Z}{\mathbb Z}
\setlength{\emergencystretch}{2em}
\title{A fully specified defining quartic in 260 variables}
\author{}\date{14 November 2026}
\begin{document}\maketitle
\begin{abstract}
The complete reference-point chart removes two rational coordinates
from the fully specified 262-variable quartic. The new polynomial has
260 total variables, 232 quadratic residuals, degree exactly four,
integer content one and 3704 monomials. Its rational input projection
agrees with the predecessor for every fixed coefficient parameter.
Consequently the same certified value $N=9682$ and the same coefficient
height $149^{38730}+1$ are retained, under exactly the same arithmetic
premises. No new prime-selection, height or pole-realization input is
introduced. The same replacement also reduces the 226-variable
existence family to 224 variables, although that is not the smallest
existence family in this collection.
\end{abstract}

\section{Statement}
All counts include the input. The source is the fully specified
262-variable family of \cite{Parent}, at the retained main revision
\texttt{fa09c9c6c97982c13265e7088d95cdb61b3e90cf}. Its quadratic-table
SHA-256 is
\begin{center}\small
\texttt{84c1130856dfda7d3a489f3077d7bc378cd8e594198dd8f3c03fe50463606a93}.
\end{center}
The coefficient parameter is $A=149^{2N+1}$, never a quantified
coordinate. Retain the exact elliptic group, normalized modular
realization, pole-realization and fixed-data prime-pattern premises
stated by that source.

\begin{theorem}\label{thm:main}
There is a completely specified family
$P_N\in\Z[X,Y_1,\ldots,Y_{259}]$, $N\ge1$, with exact degree four,
content one and 3704 nonzero coordinate monomials. For every fixed
$N$, its rational input projection is equal to that of the source
262-variable polynomial. Its height is
\[
                  H(P_N)=149^{4N+2}+1\qquad(N\ge53).
\]
Under the unchanged source arithmetic premises, every $N\ge9682$
defines $\Z$ over $\Q$. In particular $N=9682$ gives one fully
specified polynomial of height $149^{38730}+1$.
\end{theorem}

The source's effective threshold is preserved because the projection
equivalence is for every fixed $N$, not merely for sufficiently large
members with a new unspecified threshold. The algebraic equivalence
itself is unconditional. The statement that either polynomial defines
$\Z$ retains the source arithmetic inputs and their audit limitations.

\section{The exact coordinate deletion}
Work over $K=\Q(\rho)$, $\rho^2=2$. The source reference block stores
$f=1/x(H)$, $S=y(H)/x(H)$, $g=1/x(H-P)$ and $T=y(H-P)/x(H-P)$,
with the complete three-field-quadratic fixed-difference chart. It also
stores $B=f^2$ and $w=g^2$ for the local factors. The latter four
rational output coordinates have no other definition required beyond
these two squares.

In the literal source table the slots are
\[
 f=(x_{81},x_{82}),\quad S=(x_{83},x_{84}),\quad T=(x_{55},x_{56}),
 \quad g=(x_{85},x_{86}),
\]
\[
 B=x_{132}+2\rho x_{133},\qquad w=x_{87}+2\rho x_{88}.
\]
All ordinary pairs in the first line use the basis $1,\rho$.
The factor two in the second line is part of the source convention.
Delete $x_{85},x_{86}$ and their two square-definition rows 23,24.
Replace the six reduced reference rows 17--22 by the rational
components of
\[
 Q_1(f,S,B,T)=0,\qquad M(f,S,B)w+L(f,S,B,T)=0,
 \qquad Q_3(f,S,B,T)=0,
\]
using exactly the corrected coefficient table in \cite{Chart}.
The already present rows 68,69 impose $B=f^2$ and are retained.
Every other row is unchanged, apart from the final coordinate
compaction. Inspection of all occurrences, performed by the checker,
shows that deleted slots 85,86 occur only in the six reference rows
and the two deleted square definitions.

\begin{lemma}\label{lem:equivalent}
The deleted coordinates are uniquely reconstructed from every rational
solution of the new equations, and the reconstructed tuple solves the
old equations. Conversely deleting these two coordinates from any old
solution gives a new solution.
\end{lemma}
\begin{proof}
The complete chart theorem \cite[Theorem~1.1]{Chart} applies over $K$,
with no arithmetic hypothesis. It gives the actual point $H$, its
translate, and the interpretations of $f,S,B,w,T$, including all
exceptional leading-coefficient cases. In its notation,
\[
 G=\frac{((118+60\rho)S+156+69\rho/2)T-(156+69\rho/2)S+288+128\rho}
              {252-672\rho},\qquad g=G^2/f.
\]
The theorem proves $f\ne0$, so this reconstructs a unique element of
$K$, hence two rational coordinates. It also proves $w=g^2$ and all
three original fixed-difference equations. The six old reduced rows
are an invertible rational row transformation of those component
equations, so they vanish. Their coefficient convention is retained
in the source table. No other row consumes $g$, so all remaining old
rows hold. Conversely every old reference point satisfies the four
new field quadrics by the converse part of that same theorem; all
unchanged rows agree literally. The input coordinate and $N$ have
not changed.
\end{proof}

Thus the new count is $262-2=260$ coordinates and $234-2=232$ rows.
Their primitive integer sum of squares vanishes exactly when all rows
vanish, proving projection equivalence.

\section{Exact coefficient height and checks}
The full producer and independent checker are in
\begin{center}\small
\nolinkurl{computation/quartic/compressed_reference_v27}.
\end{center}
The \texttt{family260} subdirectory includes the complete quadratic
and expanded quartic tables. There are 3705 terms with powers of $A$
kept separate, and 3704 ordinary coordinate monomials. All coordinate
indices occur, degree is four, and the coefficient gcd is one.

Every coefficient profile is a fixed nonzero integer, a nonzero linear
multiple of $A$, or the single profile $A^2+1$. The largest linear
coefficient in absolute value is 44402. The full exact fixed maximum
$C_0$ is recorded as an integer in the construction certificate.
The checker proves by integer comparison
\[
 C_0<149^{214}+1,\qquad44402<149^{107}.
\]
Consequently for $A\ge149^{107}$ all constant and linear profiles
are strictly smaller in absolute value than $A^2+1$. The latter is
attained, proving the height formula for every $N\ge53$.
Positive parameter values cannot cancel any of the stated profiles,
so the monomial count is uniform. The unchanged effective source
threshold $9682$ is much larger, and is logically distinct from 53.

The independent checker imports no producer. It authenticates the
source bytes, reconstructs all changed rows from the separately
certified field-equation table, compares all unchanged rows, and
expands every coefficient by a different term-pair enumeration.
Changing a row or a quartic coefficient is rejected. The chart's
own checker separately verifies all ideal identities and every
exceptional case, including rejection of the omitted $fS$ term.
These computations are not a new audit of the source arithmetic
premises and are not represented as non-author acceptance.

\section{The corresponding 224-variable family}
The 226-variable source has the same reference slots and square
conventions. The identical deletion gives 224 variables, 198 rows,
3045 coordinate monomials and degree four. The source hash is
\begin{center}\small
\texttt{ccc845018e4d89b9d19f797bbfeca10c17ec637fbde5cb5bc08ad3bba9da105b}.
\end{center}
The proof of Lemma~\ref{lem:equivalent} is unchanged. Its coefficient
height is again $149^{4N+2}+1$ for $N\ge53$, but its sufficient
defining threshold remains ineffective, just as in that source.
The complete tables and checks are in \texttt{family224}. This is an
additional application of the exact chart equivalence, not a claim to
improve the separate 204-variable existence construction.

\begin{thebibliography}{9}
\bibitem{Parent} \emph{A fully specified defining quartic in 262 variables},
retained project paper, \nolinkurl{quartic262_20261110_v21.tex}.
\bibitem{Chart} \emph{A complete five-coordinate elliptic reference chart},
corrected project paper, \nolinkurl{compressed_reference_20261114_v27.tex}.
\bibitem{Existence} \emph{A family of defining quartics in 226 variables},
retained project paper, \nolinkurl{quartic226_20261110_v21.tex}.
\end{thebibliography}
\end{document}
